\section{5. Conclusion}


In this paper, we proposed a unified dual-mask physical model to address the persistent challenge of accurate depth estimation for non-Lambertian surfaces in light field imaging. By integrating a dual-mask mechanism with multi-directional epipolar plane image (EPI)s (EPIs), our approach effectively decouples complex specular and transparent reflections from reliable geometric cues, providing a lightweight and efficient baseline that seamlessly accommodates mixed reflectance scenarios.

The core contributions of this work are threefold:
\begin{itemize}
  \item \textbf{Unified Light Field Depth Estimation Architecture:} We constructed a dual-mask and multi-directional EPI framework that overcomes the inability of traditional models to simultaneously handle Lambertian and non-Lambertian scenes. Achieving an overall Mean Absolute Error (MAE) of 0.1795, this lightweight baseline validates the efficacy of combining multi-directional EPIs with dual masks for robust depth inference.
  \item \textbf{Theoretical Falsification of Frequency-Domain Material Classification:} We established the physical limits of angular frequency-domain analysis, making a critical "negative discovery" by rigorously proving the infeasibility of material classification under the mainstream $9 \times 9$ angular resolution. By clarifying that a minimum of $17 \times 17$ angular sampling is required, this theoretical boundary prevents blind exploration in low-resolution regimes and provides explicit guidance for light field sensor design.
  \item \textbf{Domain-Balanced Training and Architectural Limit Diagnosis:} We introduced a domain-balanced strategy to tackle long-tail distributions in multi-domain joint training, mitigating the gradient submersion of minority domains and significantly enhancing generalization in data-scarce scenarios. Furthermore, we diagnosed the inherent architectural limits of pure EPI-based methods, demonstrating that surpassing current performance bottlenecks necessitates non-EPI or hybrid architectures.
\end{itemize}
Despite the substantial advancements demonstrated by our unified model, the current validation is inherently constrained by the geometric limitations of pure EPI representations, which remain vulnerable to severe occlusions and extreme view-dependent distortions. A specific and promising future direction is to extend our dual-mask physical priors into hybrid architectures that integrate non-EPI spatial-contextual features or implicit neural representations, thereby transcending the epipolar geometry bottleneck while leveraging the theoretical angular sampling guidelines established in this study.


From a signal-processing perspective, the efficacy of the proposed dual-mask mechanism lies in its ability to perform adaptive spatial-angular filtering on Epipolar Plane Images (EPIs). In idealized Lambertian scenes, EPIs exhibit continuous, high-contrast linear structures where the slope directly corresponds to disparity. However, non-Lambertian phenomena—such as specular highlights and transparent refractions—introduce severe high-frequency angular noise and slope discontinuities. Rather than treating these artifacts as mere outliers, our Medium Mask and Angular Direction Mask collaboratively act as a learnable bandpass filter. The Medium Mask implicitly segments the optical propagation path, isolating direct illumination from complex inter-reflections, while the Angular Direction Mask dynamically re-weights the epipolar gradients. This prevents the convolutional kernels from overfitting to high-amplitude specular lobes, effectively executing a soft-thresholding operation that preserves the underlying geometric phase information without discarding critical texture cues.

The theoretical falsification of frequency-domain material classification under $9 \times 9$ angular sampling carries profound implications for next-generation light field hardware design. Our analysis demonstrates that the angular frequency spectra of complex Bidirectional Reflectance Distribution Functions (BRDFs) contain high-frequency lobes that severely alias when sampled at standard commercial resolutions. By establishing that a minimum of $17 \times 17$ angular resolution is required to satisfy the Nyquist-Shannon sampling theorem for these specific material signatures—where the maximum BRDF frequency $f_{\text{BRDF}}^{\text{max}}$ must be strictly bounded by the angular Nyquist limit:
\begin{equation}
\label{eq:eq30}
f_{\text{Nyquist}} = \frac{N_{\text{angular}}}{2} \geq f_{\text{BRDF}}^{\text{max}}
\end{equation}
—we highlight a critical bottleneck in current plenoptic camera architectures. This finding suggests that future hardware iterations must fundamentally rethink the spatial-angular resolution trade-off. Until micro-lens arrays can natively support denser angular sampling without catastrophic spatial degradation, algorithmic pipelines must either abandon pure frequency-domain material segmentation or incorporate angular super-resolution as a mandatory pre-processing stage.

Furthermore, the engineering rationale behind the EPINet4Dir V3 architecture specifically addresses the memory and computational bottlenecks that have historically hindered the deployment of light field algorithms on edge devices. By extracting multi-directional EPIs, the network captures essential 2D geometric cues while entirely bypassing the cubic computational complexity associated with full 4D light field convolutions. Constraining the parameter count to merely 754K enables end-to-end training on $384 \times 384$ full-resolution patches without exhausting the VRAM limits of embedded GPUs. This lightweight footprint, combined with the unified handling of mixed reflectances, transitions non-Lambertian depth estimation from a computationally prohibitive offline process to a viable real-time solution for autonomous navigation and robotic manipulation in unstructured environments.

Despite these advancements, the current framework exhibits notable limitations in scenarios featuring extreme occlusions coupled with highly anisotropic reflections, such as brushed metals or woven fabrics. In such regions, the EPI slope becomes highly non-linear and discontinuous, violating the local linearity assumption inherent to EPI-based disparity extraction. Additionally, while the dual-mask mechanism successfully decouples reflectance artifacts empirically, it relies on implicit feature learning rather than explicit physical rendering equations. Future research will focus on integrating differentiable rendering techniques to explicitly model anisotropic BRDFs within the mask generation process. By bridging data-driven feature extraction with rigorous physically-based rendering (PBR) constraints, we aim to further enhance the model's generalization to unseen, highly complex optical materials.

Despite the demonstrated efficacy of the unified dual-mask physical model, several limitations remain, primarily stemming from the inherent complexities of light transport and hardware constraints. First, while the medium and angular direction masks successfully decouple standard specular and transparent reflections, their generation relies on distinct angular signatures. In extreme optical scenarios involving multi-bounce refractions, complex caustics, or nested transparent layers, these signatures become highly entangled. Consequently, mask prediction may suffer from feature leakage, inadvertently suppressing valid geometric cues alongside non-Lambertian artifacts. Second, given the physical bounds of angular frequency-domain analysis at the mainstream $9 \times 9$ resolution established in this work, the network heavily compensates by leveraging spatial context. This introduces a critical vulnerability in textureless non-Lambertian regions, where the lack of spatial gradients exacerbates depth ambiguity, leading to localized smoothing or structural collapse. Furthermore, the current dual-mask formulation treats material properties as discrete categorical states rather than continuous physical parameters. This discretization limits the model's ability to handle surfaces with spatially varying, mixed micro-facet distributions, such as brushed metals or weathered plastics, where the transition between Lambertian and non-Lambertian behavior is gradual. Additionally, the current formulation assumes static illumination; the mechanism does not explicitly account for inter-reflections between adjacent non-Lambertian objects, which can introduce spurious angular variations that mimic geometric disparity.

To address these challenges, future work will explore several promising directions. A primary avenue is hardware-algorithm co-design targeting the theoretically derived $17 \times 17$ angular sampling threshold. By integrating advanced computational imaging techniques, such as coded apertures or metasurface-based light field cameras, future research will seek to acquire high-angular-resolution data, which would facilitate true frequency-domain material decoupling without relying exclusively on spatial priors. Furthermore, we plan to embed the dual-mask mechanism into differentiable rendering frameworks and neural implicit representations, such as Light Field Neural Radiance Fields (NeRF) or 3D Gaussian Splatting. By explicitly parameterizing volumetric light transport, this approach would transition the model from a discriminative masking paradigm to a generative physical inversion process, enabling the joint optimization of depth and spatially varying Bidirectional Reflectance Distribution Functions (BRDFs). To overcome the aforementioned discretization limits, future iterations will also incorporate continuous BRDF parameter estimation directly into the mask generation pipeline, allowing for sub-pixel material blending. Finally, extending the multi-directional Epipolar Plane Image (EPI) architecture to the temporal domain by formulating spatio-temporal EPIs will be investigated. Incorporating temporal consistency constraints will enhance robustness for dynamic non-Lambertian scenes, ultimately advancing light field depth estimation toward real-time, edge-deployable 3D perception in unconstrained environments.